\begin{frame}[fragile]{strlen e sizeof respondem perguntas diferentes}
\begin{lstlisting}
#include <string.h>
/* Dentro de main: */
char estado[10] = "ON";
printf("Comprimento: %zu\n", strlen(estado));
printf("Capacidade: %zu\n", sizeof estado);
\end{lstlisting}
\begin{tabularx}{\textwidth}{lXl}
\toprule
Expressão & Significado neste exemplo & Resultado \\
\midrule
\texttt{strlen(estado)} & Bytes anteriores ao terminador & 2 \\
\texttt{sizeof estado} & Bytes de todo este vetor & 10 \\
\bottomrule
\end{tabularx}
\vspace{0.25cm}
\small Os resultados usam \texttt{size\_t}, impresso com \texttt{\%zu}. \texttt{strlen} exige uma string terminada. Em UTF-8, uma letra acentuada pode ocupar mais de um byte; os comandos desta aula usam apenas ASCII.
\end{frame}
